\chapter{\texttt{straight.el} build-graph internals}
\label{app:straight-internals}
\index{straight.el!internals}\index{recipe resolution}\index{dependency tracking}\index{modification detection}\index{MELPA}

This appendix supplements Chapter~\ref{ch:straight} with a deeper
look at how \texttt{straight.el} resolves package dependencies and
decides what to rebuild.

\section{The recipe resolution algorithm}

When \texttt{straight-use-package} is called, straight resolves the
package's recipe through a priority chain:

\begin{enumerate}
  \item The literal recipe given in the \texttt{use-package}
    declaration (e.g.\ \texttt{:straight (pkg :type git :host github ...)}).
  \item A recipe registered by \texttt{straight-register-package}.
  \item The default recipe from the first recipe repository that
    knows the package (typically MELPA, then GNU ELPA).
\end{enumerate}

Recipes specify the repository URL, branch, files to include, and
build steps.  The \texttt{:files} key defaults to
\texttt{("*.el" "*.el.in" "dir" "*.info" "*.texi" "*.texinfo"
"doc/dir" "doc/*.info" ...)}.

\section{Dependency tracking}

straight.el extracts dependencies from the package's
\texttt{Package-Requires} header (for single-file packages) or
\texttt{PACKAGE-pkg.el} (for multi-file packages).  Each dependency
is recursively resolved through the same recipe chain.

If a dependency is listed in \texttt{Package-Requires} but no recipe
is found, straight signals an error at install time.  The fix is
either a \texttt{straight-register-package} call or a
\texttt{:straight (...)} recipe for the dependency's own
\texttt{use-package} declaration.

\section{Rebuild decision logic}

A package is rebuilt when straight detects that its source has
changed since the last build.  The detection strategy is
controlled by \texttt{straight-check-for-modifications}
(see Chapter~\ref{ch:straight}).

The rebuild itself:

\begin{enumerate}
  \item Copy \texttt{:files}-matched sources from
    \texttt{repos/PACKAGE/} to \texttt{build/PACKAGE/}.
  \item Byte-compile every \texttt{.el} file, producing \texttt{.elc}
    files.
  \item Generate \texttt{PACKAGE-autoloads.el} by scanning
    \texttt{;;;\#\#\#autoload} cookies.
  \item If \texttt{straight-cache-autoloads} is set, merge the
    individual autoload file into the global autoload cache.
\end{enumerate}

Native compilation (\texttt{.eln} files) is handled separately by
\texttt{comp-el-to-eln-rel-filename} and runs asynchronously after
byte-compilation completes.

\section{Symlinks versus copies}

On systems that support symlinks, straight creates symbolic links
from \texttt{build/PACKAGE/} back to \texttt{repos/PACKAGE/} for
\texttt{.el} files, so edits to source files are reflected
immediately without rebuilding.  On Windows (where symlinks
require elevated privileges), straight falls back to file copies,
and a manual \texttt{straight-rebuild-package} is required after
edits.
